\section{Purpose}
\label{sec:purpose}%

This Integration Testing Document (ITD) presents a comprehensive plan for testing the integration of components in the \textbf{Best Bike Paths (BBP)} system. The document outlines the strategy, test cases, and procedures for validating that all system components work together correctly according to the architectural design specified in the Design Document.

The main objective of this document is to provide a systematic approach to integration testing that ensures all components of the BBP system integrate seamlessly, following the Bottom-Up integration strategy defined in the Design Document. This ITD serves as a guide for developers and testers to validate component interactions, data flows, and end-to-end functionality.

\section{Scope}
\label{sec:scope}%

This document covers the integration testing strategy and test cases for the \textbf{Best Bike Paths (BBP)} system, including:

\begin{itemize}
    \item Integration testing approach and strategy
    \item Component integration plans for all system layers
    \item Integration test case specifications
    \item Integration sequence diagrams
    \item Testing tools and environment setup
\end{itemize}

The document focuses on integration testing at various levels: data layer, service layer, API Gateway, client integration, and end-to-end scenarios. It does not cover unit testing or system acceptance testing, which are documented separately.

\section{Definitions, Acronyms, Abbreviations}
\label{sec:definition_acronyms_abbreviations}%

\subsection{Definitions}
\label{subsec:definitions}%
\begin{itemize}
    \item \textbf{Integration Test:} A test that verifies the interaction between two or more components of the system.
    \item \textbf{Component Integration:} The process of combining individual components to form a larger subsystem or system.
    \item \textbf{Bottom-Up Integration:} An integration approach where lower-level components are integrated first, followed by higher-level components.
    \item \textbf{Test Case:} A specification of inputs, execution conditions, and expected outcomes for testing a particular integration scenario.
    \item \textbf{Integration Level:} A stage in the integration process representing a specific layer or grouping of components.
\end{itemize}

\subsection{Acronyms}
\label{subsec:acronyms}%
\begin{itemize}
    \item \textbf{BBP:} Best Bike Paths
    \item \textbf{ITD:} Integration Testing Document
    \item \textbf{DD:} Design Document
    \item \textbf{RASD:} Requirements Analysis \& Specification Document
    \item \textbf{API:} Application Programming Interface
    \item \textbf{DBMS:} Database Management System
    \item \textbf{UC:} Use Case
    \item \textbf{E2E:} End-to-End
\end{itemize}

\subsection{Abbreviations}
\label{subsec:abbreviations}%
\begin{itemize}
    \item \textbf{UC:} Use Case (e.g., UC1, UC2)
    \item \textbf{F:} Functional Requirement
\end{itemize}

\section{Document Structure}
\label{sec:document_structure}%
The document is organized into the following sections, each serving a distinct purpose:

\paragraph{1. Introduction:} This section presents the document's purpose, scope, and essential information required for understanding the integration testing approach. It includes definitions, acronyms, and abbreviations to facilitate comprehension.

\paragraph{2. System Implementation:} Provides a thorough overview of the implemented system, including design choices, technologies used, programming languages and frameworks, implemented features from the design document, and explicitly excluded features.

\paragraph{3. Integration Strategy and System Integration:} Describes the integration approach employed during development, following a bottom-up strategy in the backend while simultaneously developing the frontend. Covers the integration phases, synchronization between frontend and backend, and the testing methodology implemented during and after integration.

\paragraph{4. Testing Procedure:} Details the complete test structure for both backend and frontend, explaining the organization of test files, test categories (unit, integration, E2E), and the purpose of each test suite implemented in the system.

\paragraph{5. Installation Instructions:} Provides step-by-step instructions for setting up and running the system locally, including dependencies, environment configuration, and deployment procedures.

\paragraph{6. Effort Spent:} Reports the contribution of each team member to the development and documentation of this project.

\paragraph{7. References:} Provides a list of all supplementary materials and references utilized in the production of this document.
